\section{Conclusion}\label{sec:conclusion}

This paper evaluates whether a pooled linear regression remains a useful summary after the data are divided into observed subpopulations. The simulation shows that subgroup coefficient differences have more than one source. Small or uneven groups can generate substantial finite-sample disagreement even with common slopes. Shifts in predictor means can inflate raw coefficient differences and alter the precision of subgroup estimates. Slope heterogeneity produces larger persistent departures and a greater loss from using the pooled model within a subgroup.

The findings do not imply that pooled regression should be rejected whenever subgroup estimates differ. The reduced-grid interaction benchmark detects the simulated slope-heterogeneity scenarios at high rates and remains close to its nominal level when slopes are common. The proposed diagnostics answer a different question: they describe the size of subgroup departures, their relationship to covariate separation, and the within-group predictive cost of pooling. This combination supports a more limited practice: report subgroup stability when pooled coefficients are intended to describe several distinct groups, and interpret coefficient drift alongside covariate separation, conventional interaction tests, and within-group prediction performance. The empirical illustration shows how this practice can reveal a pooled model that is not equally informative across age groups.

The framework is descriptive and finite-sample. It is not a new estimator, a causal procedure, or a replacement for interaction models and stratified regression. Its value lies in making the adequacy of pooling an explicit empirical question.